% GENERATED FILE. Edit canonical lessons, then run npm run generate:books.
% canonical-source-hash: fnv1a64:e912adde6d423d08
% canonical-lessons: GU-C24-hear-savaar, GU-C24-hear-bapor, GU-C24-hear-saanj, GU-C24-hear-raat, GU-C24-day-parts-four, GU-C24-raat, GU-C24-hear-divas, GU-C24-hear-mahino, GU-C24-hear-aaj, GU-C24-aaj, GU-C24-hear-atyaare, GU-C24-time-eight, GU-R24-time-eight-r1, GU-R24-map-ten-r3

\chapter{The Day and Its Times}
\label{ch:gu-day-and-its-times}
\hlchaptermodality{28}

\begin{chapteropening}
\textbf{By the end of this chapter:} \emph{I can hear and say eight everyday Gujarati time words, and read and write two of them.}

Four parts of a day, the day, the month, today, and now enter meaning-first by ear; two of them reach the page with signs taught long before, the chapter closes on four separately scored skills, and one short return re-scores all eight a lesson later.
\end{chapteropening}

\section[savār]{\gu{સવાર} — hear morning first}

\label{lesson:GU-C24-hear-savaar}



\begin{quote}
Listen twice: \textbf{savār}. It means \textbf{morning}. Hold the sound and the meaning together; the Gujarati spelling waits for a later lesson in this chapter.
\end{quote}

\subsection*{Your turn}
Say these aloud:

\begin{itemize}
\raggedright
  \item \textquotedblleft{}savār\textquotedblright{}
  \item \textquotedblleft{}savār\textquotedblright{} when you picture early light, and nothing when you picture a dark sky
\end{itemize}

\begin{tcolorbox}[breakable,colback=teal!4,colframe=teal!35!black,title={Before you move on}]
What part of the day is \emph{savār}? (\textbf{Morning.})

Source: \href{https://dictionary.cambridge.org/dictionary/english-gujarati/morning}{Cambridge: morning}.
\end{tcolorbox}
\section[bapor]{\gu{બપોર} — hear midday first}

\label{lesson:GU-C24-hear-bapor}



\begin{quote}
Say \emph{savār} from \textbf{morning}. Now listen twice: \textbf{bapor}. It means \textbf{midday}, the middle hours that English also calls the afternoon.
\end{quote}

\subsection*{Your turn}
Say these aloud:

\begin{itemize}
\raggedright
  \item \textquotedblleft{}bapor\textquotedblright{}
  \item \textquotedblleft{}savār\textquotedblright{} or \textquotedblleft{}bapor\textquotedblright{} as you hear the English cue morning or midday
\end{itemize}

\begin{tcolorbox}[breakable,colback=teal!4,colframe=teal!35!black,title={Before you move on}]
Which of the two comes first in a day? (\emph{\textbf{Savār}} — morning.)

Source: \href{https://dictionary.cambridge.org/dictionary/english-gujarati/afternoon}{Cambridge: afternoon}.
\end{tcolorbox}
\section[sānj]{\gu{સાંજ} — hear evening first}

\label{lesson:GU-C24-hear-saanj}



\begin{quote}
Say \emph{savār}, then \emph{bapor}. Now listen twice: \textbf{sānj}. It means \textbf{evening}. The nasal in the middle is light — the tongue never quite closes.
\end{quote}

\subsection*{Your turn}
Say these aloud:

\begin{itemize}
\raggedright
  \item \textquotedblleft{}sānj\textquotedblright{}
  \item the three words in day order — \emph{savār}, \emph{bapor}, \emph{sānj}
\end{itemize}

\begin{tcolorbox}[breakable,colback=teal!4,colframe=teal!35!black,title={Before you move on}]
What does \emph{sānj} mean? (\textbf{Evening.})

Source: \href{https://dictionary.cambridge.org/dictionary/english-gujarati/evening}{Cambridge: evening}.
\end{tcolorbox}
\section[rāt]{\gu{રાત} — hear night first}

\label{lesson:GU-C24-hear-raat}



\begin{quote}
Say \emph{bapor}, then \emph{sānj}. Now listen twice: \textbf{rāt}. It means \textbf{night}. Three sounds, one beat.
\end{quote}

\subsection*{Your turn}
Say these aloud:

\begin{itemize}
\raggedright
  \item \textquotedblleft{}rāt\textquotedblright{}
  \item whichever of \emph{bapor}, \emph{sānj}, \emph{rāt} answers the English cue you hear
\end{itemize}

\begin{tcolorbox}[breakable,colback=teal!4,colframe=teal!35!black,title={Before you move on}]
What does \emph{rāt} mean? (\textbf{Night.})

Source: \href{https://dictionary.cambridge.org/dictionary/english-gujarati/night}{Cambridge: night}.
\end{tcolorbox}
\section[Practice]{The day in four parts, by ear}

\label{lesson:GU-C24-day-parts-four}



\begin{quote}
Nothing new arrives here. Say the first part of the day, then the last, before any list is read to you.
\end{quote}

\subsection*{Your turn}
Keep two scores. \textbf{Listening} goes from sound to meaning; \textbf{speaking} goes from meaning to sound. A strong score in one cannot cover a miss in the other.

\begin{tcolorbox}[breakable,colback=teal!4,colframe=teal!35!black,title={Before you move on}]
Which two of the four have you not yet seen written? (\textbf{All four} — the writing of \emph{rāt} begins in the next lesson, with signs you already know.)

Sources: \href{https://dictionary.cambridge.org/dictionary/english-gujarati/morning}{Cambridge: morning}, \href{https://dictionary.cambridge.org/dictionary/english-gujarati/afternoon}{afternoon}, \href{https://dictionary.cambridge.org/dictionary/english-gujarati/evening}{evening}, and \href{https://dictionary.cambridge.org/dictionary/english-gujarati/night}{night}.
\end{tcolorbox}
\section[Practice]{Rāt reaches the page}

\label{lesson:GU-C24-raat}



\begin{quote}
Before anything is uncovered, say the Gujarati word for \textbf{night}.
\end{quote}

\subsection*{Your turn}
Uncover \textbf{\gu{રાત}} once. Nothing in it is new: \textbf{\gu{ર}} was taught in the courtesy chapter, the \textbf{\gu{ા}} stroke in the very first one, and \textbf{\gu{ત}} beside it. Name the three parts aloud in order, then read the whole word once.

\subsection*{Writing — guided copy, model in view}
Keep the model in view and copy \textbf{\gu{રાત}} once, left to right. Repair one shape if it needs it; do not start a row.

\begin{tcolorbox}[breakable,colback=teal!4,colframe=teal!35!black,title={Before you move on}]
How many signs make \emph{rāt}, and were any of them new today? (\textbf{Three}, and \textbf{none} was new.)

Source: \href{https://dictionary.cambridge.org/dictionary/english-gujarati/night}{Cambridge: night}.
\end{tcolorbox}
\section[divas]{\gu{દિવસ} — hear day first}

\label{lesson:GU-C24-hear-divas}



\begin{quote}
Say the word for \textbf{night}, and recall the three signs that wrote it. Now listen twice: \textbf{divas}. It means \textbf{day} — the whole span that \emph{savār}, \emph{bapor}, \emph{sānj}, and \emph{rāt} divide between them.
\end{quote}

\subsection*{Your turn}
Say these aloud:

\begin{itemize}
\raggedright
  \item \textquotedblleft{}divas\textquotedblright{}
  \item \textquotedblleft{}divas\textquotedblright{} for the whole span, then \textquotedblleft{}savār\textquotedblright{} for its first part
\end{itemize}

\begin{tcolorbox}[breakable,colback=teal!4,colframe=teal!35!black,title={Before you move on}]
Is \emph{divas} a part of the day or the whole of it? (\textbf{The whole of it.})

Source: \href{https://dictionary.cambridge.org/dictionary/english-gujarati/day}{Cambridge: day}.
\end{tcolorbox}
\section[mahino]{\gu{મહિનો} — hear month}

\label{lesson:GU-C24-hear-mahino}



\begin{quote}
Say the word for \textbf{day}, and say how many signs wrote \emph{rāt}. Now listen twice: \textbf{mahino}. It means \textbf{month} — the span that holds many a \emph{divas}.
\end{quote}

\subsection*{Your turn}
Say these aloud:

\begin{itemize}
\raggedright
  \item \textquotedblleft{}mahino\textquotedblright{}
  \item \textquotedblleft{}divas\textquotedblright{} for the smaller span, then \textquotedblleft{}mahino\textquotedblright{} for the larger one
\end{itemize}

\begin{tcolorbox}[breakable,colback=teal!4,colframe=teal!35!black,title={Before you move on}]
What does \emph{mahino} mean? (\textbf{Month.})

Source: \href{https://dictionary.cambridge.org/dictionary/english-gujarati/month}{Cambridge: month}.
\end{tcolorbox}
\section[āj]{\gu{આજ} — hear today first}

\label{lesson:GU-C24-hear-aaj}



\begin{quote}
Say the word for \textbf{day}, then the word for \textbf{month}. Now listen twice: \textbf{āj}. It means \textbf{today} — this \emph{divas}, the one you are standing in. Two sounds only.
\end{quote}

\subsection*{Your turn}
Say these aloud:

\begin{itemize}
\raggedright
  \item \textquotedblleft{}āj\textquotedblright{}
  \item \textquotedblleft{}divas\textquotedblright{} for any day, then \textquotedblleft{}āj\textquotedblright{} for this one
\end{itemize}

\begin{tcolorbox}[breakable,colback=teal!4,colframe=teal!35!black,title={Before you move on}]
What does \emph{āj} mean? (\textbf{Today.})

Source: \href{https://dictionary.cambridge.org/dictionary/english-gujarati/today}{Cambridge: today}.
\end{tcolorbox}
\section[Practice]{Āj reaches the page}

\label{lesson:GU-C24-aaj}



\begin{quote}
Before anything is uncovered, say \textbf{today}, then \textbf{day}, in Gujarati.
\end{quote}

\subsection*{Your turn}
Uncover \textbf{\gu{આજ}} once. Two signs, both old friends: the standalone \textbf{\gu{આ}} that opens \emph{āvjo}, and the \textbf{\gu{જ}} of the doorway nine. Name them in order, then read the word once.

\subsection*{Writing — guided copy, model in view}
Keep the model in view and copy \textbf{\gu{આજ}} once. Repair one shape if it needs it; do not start a row.

\begin{tcolorbox}[breakable,colback=teal!4,colframe=teal!35!black,title={Before you move on}]
Which two signs write \emph{āj}, and where had you met each before? (\textbf{\gu{આ}} from \emph{āvjo}, and \textbf{\gu{જ}} from the doorway nine.)

Source: \href{https://dictionary.cambridge.org/dictionary/english-gujarati/today}{Cambridge: today}.
\end{tcolorbox}
\section[atyāre]{\gu{અત્યારે} — hear now}

\label{lesson:GU-C24-hear-atyaare}



\begin{quote}
Say \textbf{today}, name the two signs that write it, then say \textbf{month}. Now listen twice: \textbf{atyāre}. It means \textbf{now}.
\end{quote}

\subsection*{Your turn}
Three widths of time, narrowing:

\begin{itemize}
\raggedright
  \item \emph{mahino} — many days together
  \item \emph{āj} — this one whole day
  \item \emph{atyāre} — this moment inside it
\end{itemize}

Say these aloud:

\begin{itemize}
\raggedright
  \item \textquotedblleft{}atyāre\textquotedblright{}
  \item the three in narrowing order — \emph{mahino}, \emph{āj}, \emph{atyāre}
\end{itemize}

\begin{tcolorbox}[breakable,colback=teal!4,colframe=teal!35!black,title={Before you move on}]
What does \emph{atyāre} mean? (\textbf{Now.})

Source: \href{https://dictionary.cambridge.org/dictionary/english-gujarati/now}{Cambridge: now}.
\end{tcolorbox}
\section[Practice]{Payoff — eight time words in four skills}

\label{lesson:GU-C24-time-eight}



\begin{quote}
Say the first part of a day, then the word for this moment. Keep four scores below. A strong score in one skill cannot cover a miss in another.
\end{quote}

\subsection*{Your turn}
\textbf{Listening — meaning from sound.} Eight items, shuffled, once through.

\textbf{Speaking — sound from meaning.} Six items from English cues, no Gujarati model in view.

\textbf{Reading — a cold pair.} Only the two words this chapter put on the page are read; the other six stay oral until a later chapter writes them.

\subsection*{Writing — from meaning, with every model covered}
Study the two written words for five seconds, cover every model, and write them from meaning. Record writing on its own, then repair only a missed word.

\begin{tcolorbox}[breakable,colback=teal!4,colframe=teal!35!black,title={Before you move on}]
How many of the eight can you say, and how many can you write? (\textbf{Eight} by ear, \textbf{two} by hand — and the six waiting are a later chapter's work.)

Sources: \href{https://dictionary.cambridge.org/dictionary/english-gujarati/morning}{Cambridge: morning}, \href{https://dictionary.cambridge.org/dictionary/english-gujarati/day}{day}, \href{https://dictionary.cambridge.org/dictionary/english-gujarati/today}{today}, and \href{https://dictionary.cambridge.org/dictionary/english-gujarati/now}{now}.
\end{tcolorbox}
\section[Practice]{The eight time words return at R1}

\label{lesson:GU-R24-time-eight-r1}



\begin{quote}
Nothing new arrives here. Say the word for \textbf{now} before any list is read to you, then recall which four scores the payoff kept apart.
\end{quote}

\subsection*{Your turn}
Keep the same four scores the payoff kept. \textbf{Listening} runs sound to meaning; \textbf{speaking} runs meaning to sound.

\subsection*{Writing — from meaning, with every model covered}
Cover every Gujarati model and write the two words from sound alone. Record writing on its own; a strong listening score cannot cover a miss here.

\begin{tcolorbox}[breakable,colback=teal!4,colframe=teal!35!black,title={Before you move on}]
How far after the payoff did this return come? (\textbf{One lesson} — the first spacing window.)

Sources: \href{https://dictionary.cambridge.org/dictionary/english-gujarati/month}{Cambridge: month} and \href{https://dictionary.cambridge.org/dictionary/english-gujarati/now}{now}.
\end{tcolorbox}
\section[Practice]{The map words return at a third window}

\label{lesson:GU-R24-map-ten-r3}



\begin{quote}
The five newest place words have not been asked for since the chapter that taught them. Before any list is read, say \textbf{village} and \textbf{shop} in Gujarati.
\end{quote}

\subsection*{Your turn}
\textbf{Listening — meaning from sound.} Five items, shuffled, once through.

\textbf{Speaking — sound from meaning.} Three items from English cues alone.

Keep the scores apart, exactly as the payoffs that first taught these words did.

\subsection*{Writing — from meaning, with every model covered}
Cover every Gujarati model and write the two words from meaning. Record writing on its own, then repair only a missed word.

\begin{tcolorbox}[breakable,colback=teal!4,colframe=teal!35!black,title={Before you move on}]
Why does this return come here rather than beside the words that taught it? (\textbf{Because retrieval at a distance is what makes the words durable}, and a return next door would have measured nothing.)

Sources: \href{https://dictionary.cambridge.org/dictionary/english-gujarati/city}{Cambridge: city}, \href{https://dictionary.cambridge.org/dictionary/english-gujarati/school}{school}, \href{https://dictionary.cambridge.org/dictionary/english-gujarati/road}{road}, and \href{https://dictionary.cambridge.org/dictionary/english-gujarati/village}{village}.
\end{tcolorbox}
